\chapter{Simulating and fitting}
\label{ch:simular}

This chapter covers the \emph{forward model} (simulation) and the \emph{discrete
fit}, illustrated with the simplest yet most important case: an
$\alpha$-Fe calibration. Recall the distinction from
chapter~\ref{ch:intro}: \textbf{simulating} produces a spectrum from
parameters that you set; \textbf{fitting} infers those parameters from the
measured spectrum.

\section{The forward model}

Fitbauer builds the spectrum as transmission:
\[
  T(v) \;=\; b + s\,v \;-\; \sum_{k} A_k(v),
\]
where $b$ is the baseline, $s$ an optional slope and $A_k(v)$ the
\emph{absorption} of each component. The three basic discrete shapes are:

\begin{center}
\begin{tabular}{@{}lp{0.6\textwidth}@{}}
\toprule
\textbf{Component} & \textbf{Physical parameters} \\
\midrule
\textbf{Singlet} & $\delta$ (isomer shift), $\Gamma$ (linewidth),
depth. A single peak. \\[3pt]
\textbf{Doublet} & $\delta$, $\dEQ$ (quadrupole splitting), $\Gamma$,
depth. Two lines. \\[3pt]
\textbf{Sextet} & $\delta$, $\dEQ$, $\BHF$ (hyperfine field), linewidths
$\Gamma_{1,2,3}$, depth and relative intensities. Six lines. \\
\bottomrule
\end{tabular}
\end{center}

\autofig[\textwidth]{fig_sim_components}{Simulation of the three discrete
components: singlet (one peak), doublet ($\dEQ = 0{,}80$~mm/s) and sextet
($\BHF = 33$~T with intensity ratio $3:2:1$).}

\subsection{The mathematics of the model}
\label{sec:modelo-mates}

Each component is a \textbf{sum of lines} of Lorentzian shape. A line
centred at $v_0$ with linewidth $\Gamma$ (the \emph{full width at half
maximum}, FWHM —the parameter being fitted—) is
\[
  L(v;\,v_0,\Gamma) \;=\; \frac{(\Gamma/2)^2}{(v-v_0)^2 + (\Gamma/2)^2},
\]
(maximum $1$ at $v=v_0$; the half-width is $\Gamma/2$), and the absorption of a
component is
\[
  A_k(v) \;=\; d_k \sum_{j} w_j \, L(v;\,v_{0,j},\Gamma_j),
\]
with $d_k$ the \emph{depth} and $w_j$ the relative \emph{intensities} of its
lines. What distinguishes each type is \textbf{where the lines fall} and
\textbf{how many there are}:

\begin{itemize}[nosep]
  \item \textbf{Singlet}: one line at $v_0 = \delta$. The isomer
        shift $\delta$ moves the whole pattern.
  \item \textbf{Doublet}: two lines at $v_0 = \delta \pm \tfrac{1}{2}\dEQ$. The
        quadrupole $\dEQ$ is the separation between the two.
  \item \textbf{Sextet}: six lines
        \[
          v_{0,j} \;=\; \delta \;+\; \tfrac{1}{2}\,q_j\,\dEQ \;+\;
          p_j\,\frac{\BHF}{33{,}0},
        \]
        where $p_j = (\pm0{,}839,\ \pm3{,}084,\ \pm5{,}329)$~mm/s are the
        published $\alpha$-Fe line positions at 33~T (which is why the pattern
        \emph{scales} with $\BHF/33$) and $q_j=(+1,-1,-1,-1,-1,+1)$ the
        first-order quadrupole pattern. The intensities follow the
        $3:x:1:1:x:3$ ratio (§\ref{sec:textura}).
\end{itemize}

The complete spectrum is the transmission $T(v) = b + s\,v - \sum_k A_k(v)$. This
parametrisation —positions tied to $\delta$, $\dEQ$, $\BHF$ with physical
meaning— is what lets you \emph{read} the material directly from the fitted
parameters. The full derivation (treatment of the
magnetic+quadrupolar Hamiltonian, higher orders) is in appendix~\ref{ap:mates}.

\subsection{Line profile and absorber model}

In the calibration panel you choose the \menu{Line profile}:

\begin{itemize}[nosep]
  \item \textbf{Lorentzian}: the natural shape of the Mössbauer line.
  \item \textbf{Voigt}: convolution of a Lorentzian with a Gaussian of width
        $\sigma$, useful when there is instrumental broadening or a small
        unresolved spread.
\end{itemize}

The \menu{Absorber} models the thickness of the sample. The difference lies in how
the absorption $A(v)$ is converted into transmission:

\begin{itemize}[nosep]
  \item \textbf{Thin (linear)}: $T \approx 1 - A(v)$. The absorption is
        proportional to the amount of Fe; valid for low-density samples.
  \item \textbf{Thick (saturation)}: $T \propto e^{-A(v)}$ (transmission-integral
        type form). As the thickness increases, the lines \textbf{saturate} —they
        stop growing linearly and broaden—; this mode corrects for that so as not
        to bias areas or linewidths. The full correction is in
        appendix~\ref{ap:espesor}.
\end{itemize}

\begin{truco}
For dense samples or very deep lines ($>15$--$20\%$), use \textbf{Thick}:
with the thin model the relative areas (and therefore the phase proportions)
would come out biased.
\end{truco}

\subsection{Sextet intensities and texture}
\label{sec:textura}

The six lines of the sextet follow the $3:x:1:1:x:3$ area ratio. The value
$x$ (lines 2 and 5) depends on the orientation of the moments with respect to the beam:

\begin{itemize}[nosep]
  \item In \textbf{free} mode, the relative intensities are fitted.
  \item In \textbf{texture} mode, a single parameter $t$ fixes $x$ physically
        (ratio $R_{23}$, angle $\theta$), useful for oriented samples or
        thin films. The panel shows $\theta$ and $R_{23}$ derived from $t$.
\end{itemize}

The treatment of the \textbf{quadrupole} in the sextet can be first-order
(perturbative) or complete; it is selected in \menu{Fit > Fit mode}. The
detail (Kündig-type treatment of the magnetic+quadrupolar Hamiltonian) is in
appendix~\ref{ap:mates}.

\section{Simulating without fitting}

\textbf{Live simulation} is one of Fitbauer's hallmarks: instead of launching a
fit blindly, you build the model by hand and watch the spectrum change
\textbf{instantly} with each parameter. To explore the model:

\begin{enumerate}[nosep]
  \item Choose the shape of each component in its panel (\menu{Shape:}) and enable it.
  \item Set the parameters with the sliders or by typing the value;
        each control has a \emph{slider}-type bar synchronised with the numeric
        box.
  \item The plot updates \textbf{live} with the resulting model and its
        subspectra.
\end{enumerate}

\paragraph{What simulation is for.} Simulation is not just decoration; it serves
several practical functions:

\begin{itemize}[nosep]
  \item \textbf{Seeding the fit}: a non-linear fit needs a reasonable starting
        point or it falls into a false minimum. Getting the model close "by eye"
        before \menu{Fit} makes the fit faster and more reliable.
  \item \textbf{Understanding the effect of each parameter}: moving $\delta$
        shifts the whole pattern; $\dEQ$ separates the lines in pairs; $\BHF$
        opens or closes the sextet; $\Gamma$ broadens; the depth deepens. Seeing
        it live builds physical intuition.
  \item \textbf{Anticipating an experiment or comparing hypotheses}: how would a
        mixture of two phases in these proportions look? You simulate it and
        compare it with a real spectrum overlaid (§\ref{sec:comparar} of the
        previous chapter).
  \item \textbf{Teaching}: illustrating where each line comes from without needing
        data.
\end{itemize}

Each component is drawn as its own \textbf{subspectrum} (with its optional area
fill), so you can see each one's contribution to the total. The
\textbf{lock} on each parameter \emph{fixes} it: useful for simulating while
varying only one quantity, or so that the subsequent fit does not move what you
already know.

\begin{nota}
Simulating and fitting share exactly the same forward model: what you see when you
simulate is, line by line, what the optimiser will try to match. That is why a
good initial simulation is the best investment before fitting.
\end{nota}

\guicap{gui_component_panel}{Panel of a sextet component with its controls for
$\delta$, $\dEQ$, $\BHF$, linewidths and intensities, and the fixing locks.}

\section{Fitting an $\alpha$-Fe calibration}
\label{sec:ajustar-alphafe}

$\alpha$-Fe is the reference case. The procedure:

\begin{enumerate}[nosep]
  \item Load the spectrum (\menu{File > Load…}).
  \item Enable a single component of \textbf{Sextet} shape.
  \item Give approximate initial values: $\delta \approx 0$, $\dEQ \approx 0$,
        $\BHF \approx 33$~T, $\Gamma \approx 0{,}30$~mm/s.
  \item \menu{Fit > Fit}.
\end{enumerate}

The result should reproduce the six lines with $\BHF \approx 33$~T and a
reduced $\chi^2$ close to 1, as in figure~\ref{fig:fig_calibration_fit}
(chapter~\ref{ch:calibracion}). That fit is, precisely, the one marked
as calibration (§\ref{sec:definir-calib}).

\begin{truco}
If you want the fit to also refine \param{Vmax} itself against the
$\alpha$-Fe positions, enable \emph{Fit Vmax}. To avoid
degeneracy, Fitbauer automatically fixes the active $\BHF$ when the
velocity is fitted (both cannot be refined at once).
\end{truco}

\section{Fit options}
\label{sec:opciones-ajuste}

The statistical criteria are chosen in the \menu{Fit} menu / options panel:

\begin{itemize}[nosep]
  \item \menu{Likelihood}: \textbf{Gauss} ($\chi^2$ with the data $\sigma$) or
        \textbf{Poisson} (IRLS, with the model $\sigma$). Poisson is more
        appropriate with low counts.
  \item \menu{Robust loss}: \textbf{Linear} (least squares),
        \textbf{Soft L1} (smooth) or \textbf{Huber}, which downweight
        outliers.
  \item \textbf{Multistart}: repeats the fit from several starting points so as
        not to get trapped in a local minimum.
  \item \menu{Propagate calibration $\sigma$ (vmax)}: propagates the uncertainty
        of the calibrated \param{Vmax} to the fitted parameters.
\end{itemize}

\begin{nota}
When \emph{Fit Vmax} is enabled, the calibration takes part in the fit; that
is why it is worth being clear about which calibration is set (chapter~\ref{ch:calibracion}).
\end{nota}

\subsection{Constraints and physical presets}

\menu{Fit > Constraints between parameters…} lets you \textbf{tie}
parameters (for example, equalising linewidths between components or imposing
intensity relations). \menu{Fit > Physical constraint presets…} applies
common sets of constraints at a glance. Reducing the number of free
parameters stabilises the fit and avoids degeneracies.

\section{Fit quality and errors}
\label{sec:calidad}

A fit is not finished when the optimiser converges: you have to \textbf{judge
whether it is good} and \textbf{quantify the uncertainty}. For this, Fitbauer
gathers several indicators in the information panel and several error methods.

\subsection{The statistics and what they are for}

\begin{center}
\begin{tabular}{@{}lp{0.66\textwidth}@{}}
\toprule
\textbf{Indicator} & \textbf{What it measures / what it is for} \\
\midrule
$\chi^2$ & Sum of squared residuals weighted by the noise. It decreases as the
fit improves, but it depends on the number of points. \\[3pt]
\textbf{Reduced} $\chi^2$ & $\chi^2$ per degree of freedom. It is the \textbf{main
judge}: $\approx 1$ = fit compatible with the noise; $\gg 1$ = the model does not
reproduce the data (missing components, wrong centre, poor calibration);
$\ll 1$ = overfitting or overestimated noise. \\[3pt]
\textbf{dof} (degrees of freedom) & Number of points minus free parameters. Context
for the $\chi^2$. \\[3pt]
\textbf{AIC} and \textbf{BIC} & Information criteria: they reward a good fit and
\textbf{penalise adding parameters}. They serve to \textbf{compare models}
(e.g.\ two sites or three?): \emph{lower is better}. They avoid the self-deception
that "more components always fit better". \\[3pt]
\textbf{Correlations} & Fitbauer warns of pairs of highly correlated parameters
(high $|r|$): a sign of \textbf{degeneracy} (the fit cannot separate them) or that
there are too many parameters. The solution is usually to fix or tie one
(§\ref{sec:opciones-ajuste}). \\
\bottomrule
\end{tabular}
\end{center}

\begin{truco}
Do not chase a $\chi^2_r$ of exactly 1 by adding components: watch AIC/BIC and the
correlations. A model with $\chi^2_r=1{,}1$ and well-defined parameters is better
than one with $\chi^2_r=1{,}0$ full of correlated components.
\end{truco}

\subsection{Estimating uncertainties}

The statistics tell you whether the fit is good; the \textbf{errors} tell you how
much you can trust each parameter. Fitbauer offers three methods, from lower to
higher cost and reliability:

\begin{center}
\begin{tabular}{@{}p{0.34\textwidth}p{0.58\textwidth}@{}}
\toprule
\textbf{Method} & \textbf{What it does / when to use it} \\
\midrule
Covariance (default) & Errors from the covariance matrix of the fit.
Instantaneous; valid if the uncertainties are approximately Gaussian and the
correlations moderate. It is the first estimate. \\[4pt]
\menu{Bootstrap errors (MC)…} & Refits many replicas of the spectrum
resampled according to the noise (Monte Carlo) and measures the spread of each
parameter. More robust against noise and non-linearity; it costs more time. \\[4pt]
\menu{Profile-likelihood errors…} & Scans each parameter, refitting
the others, and sees how much the $\chi^2$ worsens (\emph{profile likelihood}). It
gives realistic \textbf{asymmetric} intervals; the most reliable when the
uncertainties are \textbf{not} Gaussian or there are strong correlations. \\
\bottomrule
\end{tabular}
\end{center}

\begin{nota}
Rule of thumb: use \textbf{covariance} during exploratory work and
reserve \textbf{bootstrap} or \textbf{profile likelihood} for the figures
you are going to publish, especially if the panel warns of high correlations.
\end{nota}

\section{Fitting-aid tools}

\begin{itemize}[nosep]
  \item \menu{Fit > Initialise from minima} and \menu{Auto-fit from
        minima}: they detect the line positions and propose a starting
        point. The \textbf{minima editor} also lets you \emph{curate} them
        on the plot itself: click to add a minimum, click on a
        marker to activate/deactivate it, and a counter of contributions per
        minimum.
  \item \menu{Fit > Find centre}: refines the folding centre.
  \item \menu{Fit > Fix all} / \menu{Release all}: lock or release
        all parameters at once.
  \item \menu{Fit > Batch fit…}: fits a batch of spectra by
        \emph{warm-start} (each fit starts from the previous one), ideal for
        temperature or treatment series.
  \item \menu{Fit > Undo fit}: reverts to the previous state.
  \item \menu{Fit > Fit history…}: a table with the latest finished fits
        (time, file, mode, $\chi^2_r$ and components). Each entry stores a
        \textbf{complete session snapshot}, so \param{Restore} brings that
        fit back in full (data, calibration, model and options). The history
        persists across sessions and its maximum size is configurable from
        the dialog itself.
  \item \menu{Fit > Local Ollama AI: suggest start…}: if you have a
        local model, it proposes initial values from a summary of the spectrum.
\end{itemize}

\guicap{gui_minima_editor}{Semi-manual minima editor on the plot: the
included (red) and excluded (hollow) markers are edited with the mouse and with the
side list.}

\section{Magnetic relaxation components}

In addition to singlet/doublet/sextet, Fitbauer includes components for
\textbf{superparamagnetic relaxation} and spin dynamics, useful in
nanoparticles near the blocking temperature:

\begin{itemize}[nosep]
  \item \textbf{Relaxation} and \textbf{Blume--Tjon}: line models with a
        relaxation rate $\nu$ (parameter $\log_{10}\nu$).
  \item \textbf{Néel (size)}: lognormal size distribution with anisotropy
        $K_{\text{eff}}$ and blocking temperature; \menu{Fit > Néel global
        fit (multi-T)…} fits several temperatures at once.
\end{itemize}

\begin{nota}
These models share the component interface (they are chosen in \menu{Shape:}).
Their numerical basis is developed in appendix~\ref{ap:relax} (\emph{Magnetic
relaxation models}).
\end{nota}
